\TieudeT{PHONG TỎA VẬT LÍ}
\subsubsection*{PHẦN I. Thí sinh trả lời câu hỏi từ câu 1 đến câu 18. Mỗi câu hỏi thí sinh chỉ chọn một phương án}
\setcounter{ex}{0}
%mac
%\loai{Đơn vị gia tốc}
\begin{ex}%book
Đơn vị của gia tốc là
\choice
{$N$}
{$m/s$}
{\True $m/s^2$}
{$m^2/s$}
\loigiai{
}
\haidong{7} \end{ex}
%\loai{Biểu thức}
\begin{ex}%book
Trong khoảng thời gian $t$ vật chuyển động thẳng, vận tốc thay đổi từ $v_0$ đến $v$ thì giá trị gia tốc là
\choice
{\True $a=\dfrac{v-v_0}{t}$}
{$a=\dfrac{v+v_0}{t}$}
{$a=\dfrac{v^2-v_0^2}{t}$}
{$a=\dfrac{v-v_0}{t^2}$}

\loigiai{
}
\haidong{7} \end{ex}
\begin{ex}%book
Vận tốc trong chuyển động nhanh dần đều có biểu thức nào sau đây?
\choice
{$v=$ at $-s$}
{$v=v_0-2 as$}
{$v=a-v_0. t$}
{\True $v=v_0+$ at}

\loigiai{
}
\haidong{7} \end{ex}
\begin{ex}%book
Công thức liên hệ giữa độ dịch chuyển, vận tốc và gia tốc của chuyển động thẳng nhanh dần đều nào sau đây là đúng?
\choice
{\True $v^2-v_0^2=2 a d$}
{$v^2-v_0^2=a d$}
{$v-v_0=2 a d$}
{$v_0^2-v^2=2 a d$}

\loigiai{
}
\haidong{7} \end{ex}
\begin{ex}%book
Gia tốc trong chuyển động thẳng biến đổi đều có thể được tính theo công thức nào sau đây?
\choice
{\True $a=\dfrac{\Delta v}{\Delta t}=\dfrac{v_2-v_1}{t_2-t_1}$}
{$a=\dfrac{\Delta t}{\Delta d}=\dfrac{t_2-t_1}{d_2-d_1}$}
{$a=\dfrac{\Delta t}{\Delta v}=\dfrac{t_2-t_1}{v_2-v_1}$}
{$a=\dfrac{\Delta d}{\Delta t}=\dfrac{d_2-d_1}{t_2-t_1}$}

\loigiai{
}
\haidong{7} \end{ex}
\begin{ex}%book
Công thức tính độ dịch chuyển trong chuyển động thẳng chậm dần đều nào sau đây là đúng?
\choice
{$d=d_0+v_0 t+\dfrac{1}{2} at^2$ ($a$ và $v_0$ cùng dấu)}
{\True $d=v_0 t+\dfrac{1}{2} at^2$ ($a$ và $v_0$ trái dấu)}
{$d=d_0+v_0 t+\dfrac{1}{2} at^2$ ($a$ và $v_0$ trái dấu)}
{$d=v_0 t+\dfrac{1}{2} a t^2$ ($a$ và $v_0$ cùng dấu)}	
\loigiai{
}
\haidong{7} \end{ex}
\begin{ex}%Tailieuchuan
Trong chuyển động thẳng biến đổi đều, gia tốc
\choice
{có giá trị bằng $0$}
{\True là một hằng số khác $0$}
{có giá trị biến thiên theo thời gian}
{chỉ thay đổi hướng chứ không thay đổi về độ lớn}
\loigiai{

}
\end{ex}
%\loai{Gia tốc}
\begin{ex}%book
Trong chuyển động biến đổi, gia tốc là một đại lượng
\choice
{đại số, đặc trưng cho sự biến đổi nhanh chậm của sự thay đổi vận tốc}
{đại số, đặc trưng cho tính không đổi của vận tốc}
{\True vectơ, đặc trưng cho sự biến đổi nhanh chậm của vận tốc}
{vectơ, đặc trưng cho sự biến đổi độ lớn của vận tốc}

\loigiai{
}
\haidong{7} \end{ex}
\begin{ex}%book
Đặc điểm nào sau đây cho biết vật chuyển động thẳng biến đổi đều?
\choice
{\True Gia tốc không đổi}
{Tốc độ không đổi}
{Vận tốc không đổi}
{Gia tốc tăng đều}

\loigiai{
}
\haidong{7} \end{ex}
%\loai{Chuyển động biến đổi}
\begin{ex}%book
Trong chuyển động thẳng biến đổi đều, phát biểu nào sau đây là đúng?
\choice
{\True Vận tốc tức thời có độ lớn tăng hoặc giảm đều theo thời gian}
{Gia tốc có độ lớn luôn thay đổi}
{Vectơ gia tốc luôn cùng chiều với vectơ vận tốc}
{Quãng đường đi được trong những khoảng thời gian bằng nhau thì bằng nhau}

\loigiai{
}
\haidong{7} \end{ex}
%\begin{ex}%book
%	Chuyển động nào sau đây \textbf{không} phải là chuyển động thẳng biến đổi đều?
%	\choice
%	{Viên bi lăn xuống máng nghiêng}
%	{Vật rơi từ trên cao xuống đất}
%	{\True Hòn đá bị ném theo phương nằm ngang}
%	{Quả bóng được ném lên theo phương thẳng đứng}	
%\loigiai{
%}
%\haidong{7} \end{ex}
%\loai{Chuyển động nhanh dần}
\begin{ex}%book
Phát biểu nào sau đây là  \textbf{sai}? Chuyển động thẳng nhanh dần đều có
\choice
{\True vectơ gia tốc ngược chiều với vectơ vận tốc}
{vận tốc tức thời là hàm số bậc nhất của thời gian}
{độ dịch chuyển là hàm số bậc hai của thời gian}
{gia tốc có độ lớn không đổi theo thời gian}	
\loigiai{
}
\haidong{7} 
\end{ex}
\begin{ex}%book
Một xe máy đang đứng yên, sau đó khởi động và bắt đầu tăng tốc. Nếu chọn chiều dương là chiều chuyển động của xe, nhận xét nào sau đây là đúng?
\choice
{\True $a>0,\ v>0$}
{$a<0,\ v<0$}
{$a>0,\ v<0$}
{$a<0,\ v>0$}
\loigiai{
}
\haidong{7} \end{ex}
\begin{ex}%book
Với chiều dương là chiều chuyển động, trong công thức $d=v_0 t+\dfrac{1}{2} at^2$ của chuyển động thẳng biến đổi đều. Đại lượng có thể có giá trị dương hoặc giá trị âm là
\choice
{vận tốc}
{thời gian}
{quãng đường}
{\True gia tốc}
\loigiai{
}
\haidong{7} \end{ex}
\begin{ex}%book
Trong công thức tính vận tốc của chuyển động thẳng nhanh dần đều $v=v_0+a t$ thì
\choice
{\True a luôn luôn cùng dấu với v}
{a luôn luôn ngược dấu với v}
{a luôn luôn dương}
{v luôn luôn dương}

%\loigiai{
%}
\haidong{7} \end{ex}
%\loai{Chuyển động chậm dần}
%\begin{ex}%book
%	Trong chuyển động thẳng chậm dần đều thì
%	\choice
%	{$a<0$}
%	{$v<0$}
%	{$a. v>0$}
%	{\True $a. v<0$}	
%\loigiai{
%}
%\haidong{7} \end{ex}
\begin{ex}%book
Một vật chuyển động thẳng chậm dần đều theo chiều ngược chiều dương của trục $Ox$. Dấu của gia tốc và vận tốc là
\choice
{$v>0,\ a<0$}
{\True $v<0,\ a>0$}
{$v<0,\ a<0$}
{$v>0,\ a>0$}

\loigiai{
}
\haidong{7} \end{ex}
\begin{ex}%book
Chuyển động thẳng chậm dần đều có tính chất nào sau đây?
\choice
{Độ dịch chuyển giảm đều theo thời gian}
{\True Tốc độ giảm đều theo thời gian}
{Gia tốc giảm đều theo thời gian}
{Quãng đường tăng đều theo thời gian}

\loigiai{
}
\haidong{7} 
\end{ex}
\begin{ex}%book
Trong chuyển động thẳng chậm dần đều theo chiều dương thì vận tốc luôn
\choice
{dương, gia tốc luôn dương}
{\True dương, gia tốc luôn âm}
{âm, gia tốc luôn âm}
{âm, gia tốc luôn dương}	
\loigiai{
}
\haidong{7} \end{ex}
%\loai{Chuyển động nhanh dần và chậm dần}
%\begin{ex}%book
%	Chuyển động nhanh dần đều và chậm dần đều khác nhau ở điểm căn bản là
%	\choice
%	{chuyển động nhanh dần đều luôn có vận tốc đầu khác 0}
%	{gia tốc của chuyển động nhanh dần đều âm, chậm dần đều dương}
%	{\True chuyển động chậm dần đều luôn có vận tốc đầu khác 0}
%	{gia tốc của chuyển động nhanh dần đều dương, chậm dần đều âm}	
%\loigiai{
%}
%\haidong{7} \end{ex}

\begin{ex}%book
Một vật chuyển động trên đoạn thẳng, tại một thời điểm vật có vận tốc $v$ và gia tốc $a$ là một hằng số. Chuyển động có
\choice
{gia tốc $a$ âm là chuyển động chậm dần đều}
{gia tốc $a$ dương là chuyển động nhanh dần đều}
{\True $a.v <0$ là chuyển chậm dần đều}
{vận tốc v âm là chuyển động nhanh dần đều}	
%\loigiai{
%}
\haidong{7} \end{ex}
\subsubsection*{PHẦN 2. Thí sinh trả lời từ câu 1 đến câu 4. Trong mỗi ý a), b), c), d) ở mỗi câu, thí sinh chọn đúng hoặc sai.}
\setcounter{ex}{0}
\begin{ex}%Tailieuchuan
Trong công thức tính vận tốc của chuyển động thẳng biến đổi đều $v=v_0+ at$.
\choiceTF[t]
{$v$ luôn luôn dương} 
{$a$ luôn luôn dương} 
{\True $a$ luôn luôn cùng dấu với $v$ khi vật chuyển động nhanh dần đều} 
{\True $a$ luôn luôn ngược dấu với $v$ khi vật chuyển động chậm dần đều} 
\loigiai{ 
\begin{itemchoice}
\itemch
\itemch
\itemch
\itemch
\end{itemchoice}
}
\end{ex}
\begin{ex}%Tailieuchuan
Xét một vật chuyển động thẳng biến đổi đều.
\choiceTF[t]
{\True Vận tốc tức thời của vật có độ lớn tăng đều hoặc giảm đều theo thời gian} 
{Gia tốc của vật thay đổi theo thời gian} 
{\True Vectơ gia tốc của vật có thể cùng chiều hoặc ngược chiều với vectơ vận tốc} 
{Quãng đường vật đi được trong những khoảng thời gian bằng nhau thì bằng nhau} 
\loigiai{ 
\begin{itemchoice}
\itemch
\itemch
\itemch
\itemch
\end{itemchoice}
}
\end{ex}
\begin{ex}
\impicinpar{Một học sinh lặp lại thí nghiệm kinh điển của Galileo từ thế kỷ XVI. Hình vẽ cho thấy bố trí của thí nghiệm này. Một số chuông nhỏ được treo ở những vị trí cố định trên một sợi dây song song với dốc nghiêng. Một xe lăn nhỏ được thả nhẹ từ đầu dốc, trên thân xe có gắn một que nhỏ, que này sẽ gõ làm chuông rung lên khi xe đi qua vị trí của chuông. Thí nghiệm được lặp lại nhiều lần. Khoảng cách giữa các chuông được điều chỉnh cho đến khi thời gian xe di chuyển giữa các chuông liên tiếp là như nhau. Điều này có nghĩa là tiếng chuông kêu đều đặn.}{
\includegraphics[width=4cm]{images/2025_06_04_af03f1d88c157fe9bf98g-058}}
\choiceTF[t]
{\True Chuyển động của xe là chuyển động nhanh dần đều}
{\True Kể từ cái chuông gần xe nhất, khoảng cách giữa các chuông liên tiếp luôn tăng dần}
{Đồ thị vận tốc - thời gian của xe có dạng đường thẳng không đi qua gốc tọa độ}
{\True Độ nghiêng của dốc càng lớn thì gia tốc của xe càng lớn}
\loigiai{
\begin{itemchoice}
\itemch Xe chuyển động nhanh dần đều			
\itemch Xe chuyển động nhanh dần đều nên trong các khoảng thời gian bằng nhau thì quãng đường xe đi được tăng dần
\itemch Đồ thị vận tốc của xe có dạng đường thẳng đi lên từ gốc tọa độ
\itemch Độ nghiêng của dốc càng lớn thì gia tốc của xe càng lớn		
\end{itemchoice}
}
\end{ex}
\begin{ex}
\immini[thm]{Trong một tình huống bóng đá, thủ môn xuất phát từ vạch ngang nối hai cột của khung thành, chạy thẳng lên phía trước để bắt bóng. Hình bên là đồ thị độ dịch chuyển theo thời gian của thủ môn. Điểm $A$ tương ứng với điểm xuất phát, đoạn $AB$ là một nhánh parabol, $BC$ là đoạn thẳng.}{	\begin{tikzpicture}[line cap=round, line join=round, scale=1, font=\footnotesize, >=stealth,draw=Mapcolor]
\path
(0,0)coordinate(O)
(3.5,0)coordinate(x)
(0,2.6)coordinate(y)
;

\draw [->](O)--(x);
\draw [->](O)--(y);

\node at (0,3) {$d\ (m)$};
\node at (4,0) {$t\ (s)$};

\foreach  \x  in  {1,2,...,6}
{\pgfmathtruncatemacro{\label}{\x};
\node[below]
(\label) at (.5*\x,0) {\label};
\draw[fill=white](.5*\x,0)circle(0.5pt);}

\foreach  \y  in  {1,2}
{\pgfmathtruncatemacro{\label}{10*\y};
\node[left]
(\label) at (0,\y) {\label};
\draw[fill=white](0,\y)circle(0.5pt);}

\draw[dashed] (0,1) -- (1,1) -- (1,0);

\draw[dashed]  (0,2) -- (1.5,2) -- (1.5,0);

\draw [dashed] (2,0) -- (2,2);
\draw[very thick,Mapcolor] (0,0) .. controls (0.6,0.4) and (1.15,1) .. (1.5,2);
\draw[very thick,Mapcolor] (1.5,2) -- (2,2) -- (3,0);
\node at (-0.2,0.2) {$A$};
\node at (0.8,1.2) {$B$};
\node at (1.4,2.4) {$C$};
\node at (2,2.4) {$D$};
\node at (3.2,0.2) {$E$};
\end{tikzpicture}
}
\choiceTF[t]
{Trong khoảng thời gian từ $0\ s$ đến $6\ s$, thủ môn không đổi hướng chuyển động}
{\True Thủ môn tăng tốc trong khoảng thời gian từ $0\ s$ đến $2\ s$}
{\True Tốc độ chuyển động của thủ môn từ điểm $B$ đến điểm $C$ là $10\ m/s$}
{\True Từ $4\ s$ đến $6\ s$, vận tốc chuyển động của thủ môn có giá trị $-10\ m/s$}
\loigiai{
\begin{itemchoice}
\itemch Thủ môn đổi hướng tại thời điểm $4\ s$ (điểm $B$ là đỉnh parabol) 
\itemch Đoạn $AB$ là một nhánh parabol lồi hướng lên, thể hiện gia tốc dương \itemch Đoạn $BC$ là đoạn thẳng giảm đều, độ dốc là $-10$ $\Rightarrow$ Tốc độ là $10\ m/s$
\itemch Vận tốc là $-10\ m/s$ (do độ dốc âm) trong khoảng $4\ s$ đến $6\ s$ 		\end{itemchoice}
}
\end{ex}





